\scanpage{253}
\mspage{203}

\[
\unsure{\xi|\zeta}_{i_1}*\xi_{i_2}*\cdots*\xi_{i_m}\qquad(i_1<i_2\ \unsure{\ldots|{-}}\ <i_m),
\]
en utilisant à tour de bras les \unsure{isom.|isos.} d'associativité et les $\alpha_{\xi,\xi'},\ \beta_{\xi,\xi'}$, il n'y a \unsure{pourtant|peut-être} qu'une seule façon de l'\unsure{envoyer|immerger} dans le produit $n$-uple. J'ai bien l'impression que ça doit être vrai.

(Est-ce que ça implique déjà que $\unsure{\mathcal{C}|C}$ est $1$-connexe ?)

En fait, il y a bien une unité bilatère, mais elle n'est pas dans $\mathcal{C}$, mais dans la catégorie plus \unsure{grande} \deleted{\unsure{de}}
\[
{}_{\mathrm{id}_{\mathrm{Cat}_{/S}}}\!\backslash\,{}^{\underline{\mathrm{Hom}}(\mathrm{Cat}_{/S},\mathrm{Cat}_{/S})}\ ,
\]
qui elle, comme \deleted{\unsure{toute}} catégorie
\[
{}_{\mathrm{id}_{\unsure{\mathcal{X}|\mathscr{X}}}}\!\backslash\,{}^{\underline{\mathrm{Hom}}(\mathcal{X},\mathcal{X})}\ ,
\]
est une catégorie monoïdale, i.e. un produit associatif \uline{et} unitaire. Les \deleted{\unsure{compatib.}} flèches $\alpha_{\xi,\xi'},\ \beta_{\xi,\xi'}$ sont les flèches qui résultent de la présence de cette unité. (vu que $\mathcal{C}$ est une sous-catégorie pleine), et \unsure{donc|alors} il doit devenir clair qu'on \unsure{a} les compatibilités dites
% archive sheet number: 253
